\section{Dependencias, procedimientos y reproducción}

Las construcciones utilizadas se desarrollan dentro del artículo:
el registro dodecafásico en la sección~\ref{sec:k-registro}, la clausura
posicional en la sección~\ref{sec:alpha}, la acción en la
sección~\ref{hmtII:sec:action} y el par angular en la
sección~\ref{hmtII:sec:ii-angulos}. La respuesta constitutiva de la
sección~\ref{iii:sec:constitutiva} compone esas salidas conservando
las etiquetas de hoja y orientación.

\subsection{Orden de reproducción}

La reproducción distingue el productor de coordenadas de los evaluadores
posteriores. El primero enumera los estados admisibles y aplica las reglas
de APP, TRIT y TPK. Sus salidas regionales son entradas legítimas de los
operadores posteriores del mismo desarrollo. El programa de evaluación
del vacío no se presenta como sustituto de ese productor.

El orden de lectura del paquete es: definiciones del núcleo; enumeración
y refinamiento regional; registro firmado y transformación de $K$;
clausura posicional de $\alpha$; sección de acción y retorno; par angular;
canales contractivos; operador constitutivo; momentos cíclicos y funcional
de Barbero; secciones dimensionales y comparación posterior.

Cada etapa conserva su dominio. Un archivo de cifras regionales lleva la
huella del productor y no una atribución metrológica. La tupla de doce
componentes de $K$ tiene su propietario en el capítulo del registro. Los
coeficientes de acción, la incidencia $90/120$ y las multiplicidades
$12:-1$ se leen en sus demostraciones anteriores a la evaluación. El
programa local reproduce esas composiciones con precisión declarada.

\subsection{Tres comprobaciones de funciones diferentes}

La conservación documental comprueba identidad de archivos y continuidad
de las copias. La prueba matemática establece las identidades y los límites
en sus dominios. La ejecución numérica comprueba una evaluación concreta,
su estabilidad y los residuos de las identidades. Ninguna de estas tres
comprobaciones sustituye a las otras.

El paquete adjunto incluye un manifiesto de fuentes, los textos LaTeX,
los datos regionales utilizados y los programas de comprobación. Las pruebas
de inversión, monotonía, momentos, transporte y covariancia se encuentran
en el cuerpo del artículo. Las fuentes documentales se citan por
procedencia; el lector puede comprobar los argumentos aquí desarrollados
sin utilizar un resultado numérico objetivo como generador.

La clausura posicional de $\alpha$ se demuestra en la
proposición~\ref{prop:alpha-normalizacion}. El
teorema~\ref{thm:alpha-dos-publicaciones-completas} identifica sus
prefijos con los de la representación analítica correlacionada del mismo
estado, mediante la recurrencia explícita y los cilindros compatibles.
La inversión constitutiva de
la sección~\ref{iii:sec:constitutiva} se formula en la normalización
interna indicada: los cambios de unidades se deshacen antes de aplicar
esa inversa. Estas condiciones delimitan los resultados que utilizan
las composiciones posteriores.
